% !TeX program = lualatex
% Plantilla: APUNTES formalizados completos (skill «apuntes-formalizados»)
\documentclass[11pt,a4paper,twoside]{article}

% ---------- Datos editables ------------------------------------------
\newcommand{\autor}{Autor}
\newcommand{\asignatura}{Asignatura}
\newcommand{\titulocorto}{Título corto}
\newcommand{\curso}{Salamanca, curso 2026/2027}

\usepackage[apuntes]{apuntes-oscuro}
\hypersetup{pdftitle={Título}, pdfauthor={\autor}}

% Notación propia del tema
% \newcommand{\Lop}{\mathcal{L}}

\begin{document}

% ------------------------------ Portada ------------------------------
\begin{titlepage}
\newgeometry{top=2.2cm,bottom=2.2cm,left=2.6cm,right=2.6cm}
\centering
{\Large\scshape Universidad de Salamanca\par}
\vspace{0.9cm}
{\color{fg}\rule{\linewidth}{0.4pt}}\par\vspace{0.45cm}
{\small\color{muted}\scshape Facultad de Ciencias \quad$\cdot$\quad Grado en Física\par}
\vspace{1.4cm}
{\color{acc}\scshape \asignatura\par}
\vspace{0.15cm}
{\small\color{muted}\scshape Tema N\par}
\vspace{1.0cm}
{\fontsize{24}{30}\selectfont\bfseries Título del tema\\[2pt] segunda línea\par}
\vspace{0.7cm}
{\color{fg}\rule{0.2\linewidth}{0.4pt}}\par\vspace{0.4cm}
{\small\itshape Subtema $\cdot$ Subtema $\cdot$ Subtema\par}
\vspace{0.8cm}
% Figura de portada (tikz/pgfplots, estilo «oscuro», sin solapes)
\vspace{0.5cm}
{\small $\displaystyle \text{dos o tres fórmulas clave}$\par}
\vfill
{\color{fg}\rule{\linewidth}{0.4pt}}\par\vspace{0.3cm}
\begin{minipage}[t]{0.32\linewidth}\raggedright
{\scriptsize\color{acc}\scshape Autor}\\[2pt]{\small\autor}
\end{minipage}\hfill
\begin{minipage}[t]{0.32\linewidth}\centering
{\scriptsize\color{acc}\scshape Asignatura}\\[2pt]{\small\asignatura}
\end{minipage}\hfill
\begin{minipage}[t]{0.32\linewidth}\raggedleft
{\scriptsize\color{acc}\scshape Apuntes formalizados}\\[2pt]{\small\curso}
\end{minipage}
\restoregeometry
\end{titlepage}

\setcounter{page}{2}
{\hypersetup{linkcolor=fg}\tableofcontents}

\section{Primera sección}
\subsection{Primera subsección}

\begin{definicion}[Nombre]
\begin{equation}
a\defeq b .
\end{equation}
\end{definicion}

\begin{proposicion}[Nombre]
$a=b$.
\end{proposicion}

\begin{demostracion}
\begin{enumerate}
\item \emph{Paso}: \dots
\end{enumerate}
\end{demostracion}

\begin{observacion}[Notación]
En el manuscrito \dots; con \dots{} resulta \dots
\end{observacion}

\end{document}
